\begin{recetaConImagen}{Arroz a la Leña}{90 min}{receta:arroz_lena}
    {imagenes/arrozLena.png}{Con buena leña y paciencia, este arroz tiene un sabor espectacular.}

    \begin{minipage}[t]{0.35\textwidth}
        \section*{Ingredientes}
        \begin{itemize}[leftmargin=*]
            \item Leña de romero (o, si no, de olivo)
            \item Carne (pollo y costilla de cerdo)
            \item Bastante aceite
            \item Ajos enteros sin pelar
            \item Bolitas de pimienta
            \item Pimientos
            \item Tomate
            \item Perejil
            \item Vino
            \item Caldo y/o agua
            \item Arroz
        \end{itemize}
    \end{minipage}
    \hfill
    \begin{minipage}[t]{0.6\textwidth}
        \section*{Preparación}
        \begin{enumerate}
            \item \textbf{El fuego:} Preparar la leña de romero; si no tienes, usar leña de olivo.
            \item \textbf{Sellar la carne:} En la paellera o recipiente con bastante aceite, sellar bien la carne.
            \item \textbf{El ajo y la pimienta:} Cuando la carne esté sellada, añadir los ajos sin pelar, quitándoles un poco de arriba y abajo y dándoles un golpe con la hoja del cuchillo para chafarlos ligeramente. Añadir también unas bolitas de pimienta.
            \item \textbf{El sofrito:} Incorporar después los pimientos y, por último, el tomate con bastante perejil.
            \item \textbf{Reducción:} Cuando el sofrito reduzca, verter el vino y dejar que reduzca también.
            \item \textbf{El caldo:} Añadir entonces caldo y/o agua.
            \item \textbf{El arroz:} Cuando el líquido esté hirviendo, echar el arroz.
            \item \textbf{Cocción final:} Dejar el arroz unos 20 minutos hasta que quede en su punto.
        \end{enumerate}
    \end{minipage}

    \vspace{1em}
    \begin{tcolorbox}[colback=orange!5, colframe=orange!50, title=Nota, size=small]
        La clave de esta receta está en el sabor que aporta la leña; si puedes hacerla con romero, mucho mejor.
    \end{tcolorbox}
\end{recetaConImagen}
